%!TEX root = ../template.tex
%%%%%%%%%%%%%%%%%%%%%%%%%%%%%%%%%%%%%%%%%%%%%%%%%%%%%%%%%%%%%%%%%%%
%% 1-Introduction.tex
%% NOVA thesis document file
%%
%% Chapter with introduction
%%%%%%%%%%%%%%%%%%%%%%%%%%%%%%%%%%%%%%%%%%%%%%%%%%%%%%%%%%%%%%%%%%%

\typeout{NT FILE 1-Introduction.tex}%

\chapter{Introduction}
\label{cha:introduction}

\glsresetall

\prependtographicspath{{5-Figures/Covers/}}

% epigraph configuration
\epigraphfontsize{\small\itshape}
\setlength\epigraphwidth{12.5cm}
\setlength\epigraphrule{0pt}

\section{Motivation} \label{sec.1:Motivation}

The interest in drones has grown considerably from both research and military perspectives. These systems are increasingly employed for applications such as aerial surveillance, \gls{3D} mapping, infrastructure inspection, and other data-collection tasks. However, their use is no longer limited to research, industry, or military applications, as drones have become increasingly accessible to the general public and are frequently used for recreational and hobbyist purposes. Although recreational drone use is generally harmless, it also introduces several safety and privacy concerns that must be addressed.

Airports provide a particularly relevant example of the potential risks associated with drone operations. The sighting of a drone within or in the vicinity of an airport's airspace can result in the temporary suspension of airport operations \cite{euronews2025drone,dw2025berlin,abcnews2025drone,inews2025denmark}. In Germany alone, air traffic control authorities reported 161 drone-related disruptions in 2024 and 225 in 2025 \cite{dfs2026airtraffic}. An analysis of 139 serious drone incidents near airports worldwide further concluded that such incidents are more numerous than anticipated and occur higher and farther from the airports than expected, not only in the control zones but also in the terminal manoeuvring areas \cite{pyrgies2019}. Such incidents can result in flight delays, diversions, or cancellations, potentially causing significant operational and economic consequences. An even more critical situation can arise when drones interfere with helicopter operations at hospital heliports, where disruptions may delay emergency medical transport and potentially create life-threatening circumstances.

The increasing accessibility and affordability of consumer drones may contribute to the growing frequency of such incidents. As drones become easier to acquire and operate, some users may lack sufficient knowledge of aviation regulations, restricted airspace, and the potential consequences of operating a drone in prohibited areas. Consequently, incidents may occur unintentionally, with operators potentially entering restricted airspace without being fully aware of the associated risks or regulatory requirements. Indeed, incursions of drones into terminal airspace have been attributed to the easy access of the general public to recreational drones combined with a minimal understanding of aviation regulations \cite{wang2019alert}.

For this reason many countries, including Portugal, have laws and regulations on this subject. Some of the rules applied to the piloting of drones in Portugal set forth by \gls{ANAC} include not flying over crowds of people, not flying above 120 meters, and not dropping materials from the drone. Moreover, it is necessary to go to the authorities in order to get assigned a category to be able to fly the drone legally, where each category has a slightly different sets of rules, depending on the desired purpose for the drone. Around airports, \gls{ANAC} additionally defines prohibited and altitude-limited zones. A geographic rule, however, is of little use if the operator is unaware of it.

Technical measures have also been introduced to prevent such incursions. Manufacturers enforce no-fly zones around airports through geofencing implemented in the drone's firmware \cite{schiller2023}, and keep-out geofences are considered a building block of future unmanned traffic management \cite{kim2022geofence}. Geofencing, however, relies on the intruder's own flight controller: it is absent from drones that are not equipped with it, and firmware flaws allow it to be disabled or bypassed \cite{schiller2023}. Once a drone has entered restricted airspace, the response is therefore left to external counter-drone systems. These systems combine detection, based on radar, radio-frequency, acoustic and optical sensors, with mitigation, which comprises the jamming and spoofing of the command and navigation links, physical capture with nets, and destruction \cite{wang2021cuas,kang2020,lykou2020}. Near an airport, each of these options has important limitations. Jamming risks interfering with instrument landing systems, approach radars and radio communications, and a jammed drone may fall uncontrolled to the ground, so kinetic and electronic mitigation is often unsuitable at airports because of the collateral hazards to manned aviation \cite{lykou2020}. Teams of defensive drones have been proposed to chase an intruder at close range until it is neutralized \cite{castrillo2022}, and interceptor drones capture it with nets carried by one \cite{pliska2024} or several \cite{wang2025net,rothe2025} \glspl{UAV}. All of these approaches, however, end in the neutralization of the intruder, and none of them addresses the operator who is simply unaware of the restriction and who would turn back if warned.

This motivates the question of how a target entering restricted airspace can be safely surrounded while minimizing the risk and operational impact of the intervention while  also indicating to the target that it's entering restricted areal space. This thesis focuses on methods of guaranteeing encirclement of an evader without the use of destructive methods, that is, on a non-destructive response in which a swarm of chaser drones forms a visible barrier ahead of the intruder and, if it proceeds, surrounds it.

\section{Problem Statement} \label{sec.1:Problem Statement}

The guaranteed capture of an agent such as a drone, ship, or aircraft has been a subject of research for over a century, dating back to early work in game theory \cite{game_theory_1925}. Closely related is the problem of determining desirable formations for encircling a target, whether fixed or moving.

Despite extensive research, most existing encirclement strategies are formulated in \gls{2D} environments. In particular, current multi-agent encirclement methods predominantly focus on planar pursuit scenarios, with the majority of implementations targeting \gls{UAV}. Representative examples include circular formation control and encirclement strategies in \gls{2D} space \cite{Litimein2021, Jin2018, Kim2007, 2D_encirclement_guaranteed_sectors, Lin2004}. Existing \gls{3D} approaches, on the other hand, commonly divide the surrounding space into predefined sections and seek to guarantee that one or more chasers occupy specific regions around the target, keep the agents on a circle around a target that moves in \gls{3D} \cite{franchi2016}, or spread them over a sphere of fixed radius \cite{huang2025,bo2023}.

Furthermore, the term 'encirclement' is not defined uniformly across the literature, with different studies adopting different criteria to describe both the underlying problem and the resulting formation. Some approaches define encirclement as the target being contained within the convex region formed by the chasers \cite{mycielskiTheoriesPursuitEvasion1988,Wang_2021,Patra2025}. Other formulations consider the overlap between the target's reachability region and the regions reachable by the chasers, while others focus primarily on positioning agents around the target in a circular or spherical formation \cite{wangAchievingSafeMinimum2024,huang2025}. These definitions correspond to different objectives and, consequently, lead to different formation requirements.


In this thesis, the latter interpretation is of primary interest. The objective is not necessarily to guarantee the physical containment or capture of the target, but rather to establish a formation that provides a clear and visually observable indication that the target has entered, or is approaching, restricted airspace. Although complete encirclement remains a desirable outcome, the formation is primarily intended to function as a visible barrier or deterrent to the target and, consequently, to its operator. A formation of agents around or in front of the target is therefore considered a more appropriate objective than containment within a convex region alone. The problem considered in this thesis therefore emphasizes the formation of a visually observable spatial barrier that communicates both the presence of the chasers and the restricted nature of the surrounding airspace. 

The geometry of the restricted airspace adds further requirements. A deterrent should act before the intruder reaches the prohibited zone, so the swarm should first stand as a barrier at the boundary of the restricted airspace and only close around the target once it crosses that boundary. In the existing \gls{3D} approaches, however, the shape of the formation is fixed in advance and does not depend on where the target is relative to the airspace it threatens \cite{franchi2016,huang2025,bo2023}.

This emphasis on \gls{2D} pursuit and on containment-oriented definitions of encirclement limits the applicability of existing approaches to the considered scenario. In a fully \gls{3D} environment, altitude, spatial constraints, and maneuverability play a critical role in determining the formation and coordination of multiple agents. Addressing these limitations motivates the development of an encirclement strategy that explicitly accounts for \gls{3D} motion and coordination.

\section{Proposed solution and contributions} \label{sec.1:Proposed solution}

In this thesis, the focus is on achieving encirclement in a \gls{3D} environment. Specifically, the work addresses, the creation of positions around a target, their navigation and the generation of collision-aware trajectories during this process. The objective is to ensure safe and coordinated motion while achieving and maintaining an encirclement configuration in \gls{3D} space, with the intended operational effect of discouraging entry into, or encouraging departure from, a designated restricted aerial space.

The proposed solution can be separated into three main categories, centralized, decentralized formations creating a two-sphere shape around the target and a generalized version for a desired formation of any arbitrary well-defined surface mesh. The generalized formation shape method is used for not only ensuring encirclement around a target if the same is in restricted aerial space but also to create a physical barrier discouraging the entrance of the same into the restricted aerial space. The agents were chosen to be as equally spaced as possible in the formation independently of the formation shape, this is due to the unknown behavior of the target guaranteeing independently of the desired motion of the same that all areas are equally covered by the drone formation. 

The scope of this thesis is limited to the creation, navigation, and maintenance of an encirclement formation around a target in a \gls{3D} environment. The proposed methods are not intended to address operators with malicious intent who deliberately disregard the formation, nor targets whose dynamics or autonomy prevent them from being safely influenced or redirected. Furthermore, this work does not consider physical capture, forced interception, or collision-based methods for stopping or immobilizing a target. If a target remains within the restricted aerial space after encirclement, any subsequent intervention, decision-making, or enforcement action falls outside the scope of this thesis and is left to the relevant authorities. Accordingly, the thesis focuses on the generation and maintenance of safe, collision-aware formations and trajectories. The effectiveness of such formations in influencing operator behavior is not evaluated in this work, as well as the intent classification, the assessment of intent, or the legal and operational procedures that may follow encirclement.

This thesis contributes a framework for the generation, navigation, and maintenance of collision-aware \gls{3D} encirclement formations for aerial agents. The proposed framework addresses both centralized and decentralized formation generation and incorporates trajectory planning subject to realistic dynamic and collision constraints. The main contributions of this work are:

\begin{itemize}
\item \textbf{Comparison of point-distribution methods on the two-sphere:} a systematic comparison of methods, both directly computed or iterative, for distributing agents approximately uniformly over a spherical surface, with the objective of obtaining well-spaced encirclement configurations and uniform spatial coverage around a target.

\item \textbf{Comparison of path-planning and control approaches:} an evaluation and comparison of different path-planning approaches for navigating the agents toward the desired formation, considering their ability to achieve coordinated and collision-aware motion.

\item \textbf{Lyapunov-based target-encirclement controller:} development of a Lyapunov-based control strategy for guiding agents toward and maintaining an encirclement configuration around a moving target, providing a stability-based formulation for the formation-control problem. This contribution was published in~\cite{cavalheiroHierarchicalNonlinearControl2026}.

\item \textbf{Decentralized formation-maintenance architecture:} development of a decentralized architecture in which agents coordinate their relative positions to maintain a well-formed and approximately uniformly distributed mesh around the target, reducing the reliance on a centralized formation controller.

\item \textbf{Generalized formation-shape generation:} development of a generalized method for generating formations with arbitrary desired shapes around a target. This extends the spherical encirclement formulation to configurations that can be used both to surround a target within a restricted aerial space and to form a spatial barrier intended to discourage entry into the restricted region.

\item \textbf{\gls{MPC}-based collision-aware trajectory planning:} formulation of a \gls{MPC}-based trajectory-planning approach capable of incorporating realistic constraints on velocity and acceleration, while explicitly accounting for collision avoidance between agents and with the surrounding environment.

\item \textbf{Integration of formation generation, control, and trajectory planning:} integration of the proposed formation-generation, formation-control, and trajectory-planning methods into a unified framework for the safe creation and maintenance of \gls{3D} encirclement formations.

\end{itemize}

\section{Thesis Organization} \label{sec.1:Report organization}

The rest of this thesis is separated into multiple chapters, each chapter focuses on a single subject being the three development chapters made to build on each previous one, each chapter talks about the following subjects:

Chapter~\ref{cha:State of the art} (State of the art) refers some contextualization about the previous available works on encirclement, agent spacing, collision avoidance, trajectory planning, system model and internal drone controllers, referring the most common and well known methods.

Chapter~\ref{cha:Background} (Background) focuses on a more mathematical approach to some themes that are of great relevance for the complete understanding of the work developed. It was considered a necessary section due to its distinct theming compared with the state of the art.

Chapter~\ref{cha:Centralized Development} (Centralized Development) presents methods for achieving an almost evenly spaced distributions of points on the two-sphere, as well as the trajectory planning between current agent position and desired position on the two-sphere. These methods are used to define desired agent configurations around the target. It is also shown different methods to create a collision-aware navigation, finalizing with simulation results for a 3-agent and 31-agent swarm.

Chapter~\ref{cha:Decentralized Development} (Decentralized Development) changes the objective displacement to iterative methods for achieving almost evenly spaced points on the two-sphere, showing results for the two implemented methods.

Chapter~\ref{cha:Generalized Development} (Generalized Development) takes the encirclement problem and generalizes it having the swarm formation take shape of arbitrary well-defined surface meshes, not only the two-sphere, this leads to the possibility of not only encircling the target but prevent its entrance to restricted air space by creating a physical barrier. The restricted airspace is described by real polygonal boundaries, and an airspace-aware wall frame places a morphing formation mesh (a disc on a virtual wall, a cone, a sphere around the target and a disc on the border of the prohibited zone) relative to them. The agents estimate the target by consensus, obtain their goals from a geodesic, optionally velocity-weighted, centroidal Voronoi tessellation on the mesh, and plan their trajectories with a distributed Bernstein-polynomial \gls{MPC} solved by \gls{ADMM}, which enforces velocity, acceleration, inter-agent separation and airspace constraints through projections. The chapter ends with the Lisbon airport scenario.

Chapter~\ref{cha:Conclusions and Future Work} (Conclusions and Future Work) takes into account all previous chapters and reflects upon the results, use cases and future advances in this field. 

Appendix~\ref{cha:Appendix} presents additional methods and results that were investigated throughout the development of this thesis but were not selected for the final proposed solution. These include methods evaluated for both direct and iterative point distribution on the two-sphere, alternative approaches to agent spacing that did not achieve the desired formation properties, and additional path-planning controllers. The appendix also provides the specifications of the computer system used for the simulations and presents the comparative results obtained across the evaluated methods.

\section{Notation}

This section defines the notation used throughout the remainder of this thesis.

The notation $\mathbf{a} \cdot \mathbf{b}$ denotes the inner product between $\mathbf{a}$ and $\mathbf{b}$. While $S(\cdot)$ is the cross-product skew-symmetric matrix satisfying $S(\mathbf{a}) \mathbf{b} = \mathbf{a} \times \mathbf{b}$, for $\mathbf{a},\mathbf{b}\in\mathbb{R}^3$ and its inverse mapping $S^{-1}(\cdot)$ defined in Equation~\eqref{eq.4:skew_matrix}.

\begin{equation}
    S(\mathbf{r})
=
\begin{bmatrix}
0 & -r_3 & r_2 \\
r_3 & 0 & -r_1 \\
-r_2 & r_1 & 0
\end{bmatrix},
\qquad
S^{-1}\left(S(\mathbf{r})\right) = \mathbf{r}, \mathbf{r} \in \mathbb{R}^3
\label{eq.4:skew_matrix}
\end{equation}

Throughout this thesis, regular lowercase letters denote scalars $'x'$, calligraphic uppercase letters denote sets $'\mathcal{C}'$, bold lowercase letters denote vectors $'\mathbf{x}'$, and bold uppercase letters denote matrices $'\mathbf{X}'$. The Euclidean norm is $\|\cdot\|$, $\mathbf{1}$ is a vector of ones, $\mathbf{I}_k$ the $k\times k$ identity matrix, $\otimes$ the Kronecker product, $\mathbf{e}_3=[0\;0\;1]^{\mathsf T}$ the vertical unit vector, and $\operatorname{clip}(x,l,h)=\min(\max(x,l),h)$. $\Pi^{r,\mathbf{p}}(\cdot)$ denotes the projection onto the sphere of radius $r$ centered at $\mathbf{p}$ (Chapter~\ref{cha:Centralized Development}) and $\gls{proj}(\cdot)$ the Euclidean projection onto a set $\mathcal{C}$ (Section~\ref{sec.3:ADMM}). Discrete time instants are indexed by $[k]$, and the samples of a prediction horizon by $[m]$. The symbols used are listed in the List of Symbols, and the same symbol denotes the same quantity in every chapter.

